\documentclass[11pt,a4paper]{article}
\usepackage[T1]{fontenc}
\usepackage[utf8]{inputenc}
\usepackage{lmodern,amsmath,amssymb}
\usepackage[margin=25mm]{geometry}
\usepackage[hidelinks]{hyperref}
\hypersetup{pdftitle={Every mark with body-independent noise is a momentum-transfer mark},pdfauthor={navstokgap research working notes}}
\newcommand{\tightlist}{\setlength{\itemsep}{0pt}\setlength{\parskip}{0pt}}
\setlength{\emergencystretch}{3em}
\setcounter{secnumdepth}{0}
\begin{document}
\hypertarget{every-mark-with-body-independent-noise-is-a-momentum-transfer-mark}{%
\section{Every mark with body-independent noise is a momentum-transfer
mark}\label{every-mark-with-body-independent-noise-is-a-momentum-transfer-mark}}

Theorem A of the \href{mark-cost-and-statistical-floor.md}{mark-cost
note} was proved for the von Neumann coupling \(\lambda\hat y\hat P_A\).
Among marks whose noise does not depend on the body, that class is all
there is. For an arbitrary unitary coupling of the body to an apparatus
and an arbitrary pointer, suppose the reading error, the delivered
impulse and the position disturbance are operators of the apparatus
alone, independent of the body's state. Then:

\begin{enumerate}
\def\labelenumi{\arabic{enumi}.}
\tightlist
\item
  the error \(\hat N\) and the impulse \(\hat D\) are canonically
  conjugate, \([\hat N,\hat D]=-i\hbar\), and the position disturbance
  \(\hat X\) commutes with both (Proposition 1);
\item
  if the mark leaves the position undisturbed, \(\hat X=0\), the
  coupling is \(e^{-i\hat y\otimes\hat B/\hbar}\) followed by an
  apparatus unitary, with the pointer canonically conjugate to
  \(\hat B\): Theorem A's form exactly, with \(\hat B\) in place of
  \(\lambda\hat P_A\) and extra apparatus degrees of freedom allowed
  (Proposition 2);
\item
  the recoil bound \(s\sum_j\Delta_j\ge8(1-2\epsilon)\hbar\) of the
  \href{record-costs-recoil.md}{recoil note} holds for every protocol of
  such marks, including those that disturb position (Proposition 3).
\end{enumerate}

This is option 1 of the plan the user set on 2026-09-22 for extending
Theorem A. The class is Ozawa's measurements with \emph{independent
intervention}, for which he proves the Heisenberg product
\(\epsilon\eta\ge\hbar/2\)
(\href{https://doi.org/10.1103/PhysRevA.67.042105}{PRA \textbf{67},
042105, 2003}; \href{https://doi.org/10.1016/j.aop.2003.12.012}{Ann.
Phys. \textbf{311}, 350, 2004}; metadata); the joint-measurement
analogue is Arthurs and Goodman
(\href{https://doi.org/10.1103/PhysRevLett.60.2447}{PRL \textbf{60},
2447, 1988}, metadata). Proposition 2 is a unitary-dilation form of the
classification of translation-covariant instruments on the line (Davies,
\emph{Quantum Theory of Open Systems}, 1976, metadata), stated here in
the language of marks. What is new is the transfer to the Galileo
comparison, Proposition 3. The class needs a pointer with continuous,
unbounded spectrum, since the pointer and the impulse form a canonical
pair; marks with a discrete or bounded pointer lie outside it and are
covered by the \href{record-costs-disturbance.md}{disturbance note}.
Exploratory; no ledger promotion.

\hypertarget{the-identity}{%
\subsection{1. The identity}\label{the-identity}}

A mark is a unitary \(U\) on body \(\otimes\) apparatus, applied
impulsively, followed by reading a pointer \(\hat R\) of the apparatus.
In the Heisenberg picture define

\[\hat N=U^\dagger\hat RU-\hat y,\qquad \hat D=U^\dagger\hat pU-\hat p,\qquad
\hat X=U^\dagger\hat yU-\hat y,\]

the error of the reading as an estimate of the position before the mark,
the delivered impulse, and the displacement of the position.

\textbf{Proposition 1.} Always
\([\hat N,\hat D]=-i\hbar-[\hat y,\hat D]-[\hat N,\hat p]\). If
\(\hat N\), \(\hat D\) and \(\hat X\) commute with \(\hat y\) and
\(\hat p\), then

\[[\hat N,\hat D]=-i\hbar,\qquad[\hat X,\hat N]=[\hat X,\hat D]=0,\]

so \(\delta^2\Delta^2-c^2\ge\hbar^2/4\) with \(c\) the symmetrized
covariance of error and impulse, for every joint state.

\emph{Proof.} \(\hat R\) acts on the apparatus and \(\hat p\) on the
body, so
\([U^\dagger\hat RU,U^\dagger\hat pU]=U^\dagger[\hat R,\hat p]U=0\),
that is \([\hat y+\hat N,\hat p+\hat D]=0\); expanding with
\([\hat y,\hat p]=i\hbar\) gives the identity. Under the hypothesis the
cross terms vanish. The same argument with
\([U^\dagger\hat RU,U^\dagger\hat yU]=0\) gives \([\hat N,\hat X]=0\),
and the preservation of the canonical relation,
\([\hat y+\hat X,\hat p+\hat D]=i\hbar\), gives \([\hat X,\hat D]=0\).
The last statement is the Robertson--Schrödinger inequality. \(\square\)

Commutation with \(\hat y\) and \(\hat p\) is meant in the strong sense,
with their spectral projections, equivalently with the Weyl operators.
Operators commuting in that sense with \(\hat y\) and \(\hat p\) commute
with every body operator, so the hypothesis says exactly that
\(\hat N\), \(\hat D\), \(\hat X\) are apparatus operators.

\hypertarget{the-class-is-the-momentum-transfer-class}{%
\subsection{2. The class is the momentum-transfer
class}\label{the-class-is-the-momentum-transfer-class}}

\textbf{Proposition 2.} Suppose \(\hat X=0\) and \(\hat N\), \(\hat D\)
are apparatus operators. Then
\(U=(\mathbf 1\otimes U_0)\,e^{-i\hat y\otimes\hat B/\hbar}\) for an
apparatus unitary \(U_0\) and a self-adjoint apparatus operator
\(\hat B\), the delivered impulse is \(\hat D=-\hat B\), and
\(\hat R_0=U_0^\dagger\hat RU_0\) satisfies
\([\hat R_0,\hat B]=i\hbar\), with \(\hat N=\hat R_0\).

\emph{Proof.} \(U^\dagger\hat yU=\hat y\) means \(U\) commutes with
\(\hat y\), so in the position representation \(U=\int^\oplus U(y)\,dy\)
with apparatus unitaries \(U(y)\), measurable in \(y\). Exponentiate the
two hypotheses. Since \(\hat D\) commutes with \(\hat p\),
\(U^\dagger e^{ia\hat p/\hbar}U=e^{ia(\hat p+\hat D)/\hbar}=e^{ia\hat p/\hbar}e^{ia\hat D/\hbar}\),
and in the decomposition, where \(e^{ia\hat p/\hbar}\) translates \(y\)
by \(a\), this reads \(U(y)^\dagger U(y+a)=e^{ia\hat D/\hbar}\) for
almost every \(y\) and every \(a\). The measurable solutions are
\(U(y)=U_0e^{iy\hat D/\hbar}\), so
\(U=(\mathbf 1\otimes U_0)e^{-i\hat y\otimes\hat B/\hbar}\) with
\(\hat B=-\hat D\), and no differentiability is needed. For the pointer,
\(\hat y\) commutes with \(\hat N\), so
\(U^\dagger e^{ib\hat R}U=e^{ib\hat y}e^{ib\hat N}\); in the
decomposition, with \(\hat R_0=U_0^\dagger\hat RU_0\), this reads
\(e^{iy\hat B/\hbar}e^{ib\hat R_0}e^{-iy\hat B/\hbar}=e^{iby}e^{ib\hat N}\)
for every \(y\) and \(b\). At \(y=0\) it gives \(\hat N=\hat R_0\), and
for every \(y\) it is the Weyl relation between \(\hat B\) and
\(\hat R_0\), which contains \([\hat R_0,\hat B]=i\hbar\). \(\square\)

By the Stone--von Neumann theorem, which applies to the Weyl relation
just obtained, the pair \((\hat R_0,\hat B)\) is a canonical pair on one
factor of the apparatus, \(L^2(\mathbb R)\otimes\mathcal K\), and
\(\mathcal K\) carries whatever else the apparatus holds. So a mark that
leaves the position alone and whose noise is independent of the body is
a von Neumann mark: it couples the body's position to one canonical
momentum of the apparatus and reads the conjugate coordinate. Theorem A
and the recoil theorem apply to it verbatim, with \(\lambda\hat P_A\)
replaced by \(\hat B\) and the probe's state replaced by its reduced
state on the canonical factor, entangled with \(\mathcal K\) or not;
Theorems B and C apply with their Gaussian hypothesis placed on that
reduced state.

\hypertarget{marks-that-also-displace-the-body}{%
\subsection{3. Marks that also displace the
body}\label{marks-that-also-displace-the-body}}

A mark with \(\hat X\ne0\) adds the jump \(\hat X_j\) to the body's
position at \(t_j\), and every later reading carries it. Proposition 1
makes \(\hat X_j\) commute with the pair \((\hat N_j,\hat D_j)\), so it
acts on the multiplicity factor \(\mathcal K\). An apparatus state that
entangles \(\mathcal K\) with the canonical factor can correlate
\(\hat X_j\) with the error and the impulse, so the Gaussian floors of
the mark-cost note need that correlation bounded too. The recoil bound
needs nothing more.

\textbf{Proposition 3.} Theorem R of the recoil note holds for every
protocol of marks whose \(\hat N_j\), \(\hat D_j\), \(\hat X_j\) are
apparatus operators, with \(\Delta_j\) the spread of \(\hat D_j\).

\emph{Proof.} Lemma 2 of the recoil note needs a unitary on apparatus
\(j\) that translates \(\hat N_j\) by \(c_j\) and leaves everything else
in the record unchanged. The record depends on apparatus \(j\) through
\(\hat N_j\), \(\hat D_j\) and \(\hat X_j\) only. The unitary
\(e^{ic_j\hat D_j/\hbar}\), applied to the initial state of apparatus
\(j\), translates \(\hat N_j\) by \(c_j\) because
\([\hat N_j,\hat D_j]=-i\hbar\) (for the von Neumann mark it is the
translation of \(\hat Q_j\) by \(\lambda_jc_j\)), and fixes \(\hat D_j\)
and \(\hat X_j\), which commute with \(\hat D_j\). Mandelstam--Tamm
bounds its effect on the state by \(|c_j|\Delta_{\hat D_j}/\hbar\), and
the rest of the proof of Theorem R is unchanged. An apparatus state
entangled across marks is covered too: the \(\hat D_j\) commute, so the
translations combine into one unitary generated by
\(\sum_jc_j\hat D_j\), and the seminorm property
\(\Delta(\sum_jc_j\hat D_j)\le\sum_j|c_j|\Delta\hat D_j\) gives the same
bound. \(\square\)

\hypertarget{consequence-for-state}{%
\subsection{4. Consequence for STATE}\label{consequence-for-state}}

Option 1 is settled. Body-independent noise forces the canonical pair,
the position-undisturbed marks with such noise are exactly the von
Neumann marks, and the recoil bound
\(s\sum_j\Delta_j\ge8(1-2\epsilon)\hbar\) covers every mark with
body-independent noise, position-displacing ones included. What remains
outside is noise that depends on the body, where \([\hat y,\hat D]\) or
\([\hat N,\hat p]\) survives in Proposition 1; that is Ozawa's setting
and option 2.
\end{document}
